\paragraph{A criterion for four cases of Chowla's conjecture at once.}
Take $q=8$ and let $d\equiv7\pmod8$ be prime. Write
$\chi_D(n)=(D/n)$ for the primitive quadratic character of fundamental
discriminant $D$. Then
\[
 B_d:=\lim_{X\to\infty}\frac1{\log X}
       \int_1^X R_{8,d}(x)\frac{dx}{\sqrt{x}}
   =-\sum_{D\in\{-d,4d,-8d,8d\}}
       \operatorname{ord}_{s=1/2} L(s,\chi_D)\ \in -2\mathbb Z_{\ge0}.
\]
Indeed \eqref{eq:packet-remainder} identifies the mean with $b_0$.
The untwisted functions are $\zeta$ and the characters of discriminants
$-4,8,-8$. The alternating eta series gives $\zeta(1/2)<0$, and the
alternating beta series gives $L(1/2,\chi_{-4})>0$. For $8$ and $-8$
the signs on each block of four successive odd integers are respectively
$(+,-,-,+)$ and $(+,+,-,-)$. The sums of these blocks are positive because
$u^{-1/2}$ is strictly decreasing and strictly convex. The twist is
$\chi_{-d}$ and the four resulting discriminants are the ones displayed.
Finally each completed quadratic $L$-function has root number $+1$, so
each central order is even. Thus $B_d>-2$ is equivalent to
simultaneous central nonvanishing of all four functions. Proving this
strict lower bound for every such $d$ would establish this subfamily of
Chowla's conjecture. The one-sided estimate at exponent $-1/2+\varepsilon$
allows every fixed value of $B_d$ and supplies no such strict bound.
For the conjecture and classical nonvanishing results, see K.~Soundararajan,
\emph{Nonvanishing of quadratic Dirichlet $L$-functions at $s=1/2$},
Ann. of Math. 152 (2000), 447--488,
\href{https://doi.org/10.2307/2261390}{doi:10.2307/2261390}.
